\chapter*{Abstract}

Organizations across retail, energy, and mobility increasingly rely on
incentives to influence consumer behavior, yet most incentive systems remain
static, costly, and behaviorally uninformed. Traditional approaches—such as
blanket discounts, loyalty points, or heuristic promotional rules—fail to
account for how individuals form habits, respond to incentives, and evolve over
time. As a result, firms routinely overspend on promotions while missing
opportunities to shape long-term behavioral trajectories.

This thesis presents the Behavioral Engine for Demand (BED), a closed-loop,
state-based framework for timing-optimized incentive design. BED models
individual behavior using three latent variables—recency, habit, and
responsiveness—and formulates incentive allocation as a Constrained Markov
Decision Process (CMDP) that respects budget and fairness constraints. The
framework integrates behavioral modeling, hazard-based event prediction,
constrained optimization, simulation, and prototype implementation into a
unified research architecture.

BED v3 introduces four major contributions. First, it formalizes a behavioral
state model that captures how incentives influence habit formation and
responsiveness over time. Second, it develops a CMDP formulation that allocates
incentives efficiently under operational constraints. Third, it provides a
complete simulation environment that evaluates policy performance using
synthetic populations and behavioral dynamics. Fourth, it implements a
serverless QR-to-cloud prototype that demonstrates how BED can be deployed
safely and scalably in real-world settings.

Together, these components establish BED as a practical and theoretically
grounded approach to demand shaping. The framework enables organizations to
transition from static, intuition-driven incentives to adaptive, data-driven,
and behaviorally informed policies that optimize long-term outcomes.
